%*************************************
% ภาคผนวก จ.
%*************************************
\appendixChapterTitle{ค}{การ Deploy และการตั้งค่าสำหรับ Production}
\section{การ Deploy และการตั้งค่าสำหรับ Production}
\hspace*{1.3em}
การนำระบบ Cloud-Based Platform ไปใช้งานจริง (Production) ต้องมีการตั้งค่าและการเตรียมการที่แตกต่างจากการใช้งานในโหมด Development โดยมีรายละเอียดดังนี้

\vspace{0.5em}
\textbf{การเตรียมสภาพแวดล้อม Production}
\vspace{0.5em}

\begin{mycustomenum2}
    \item \textbf{ข้อกำหนดของ Server:}
    \begin{mycustomenum2}
        \item CPU: อย่างน้อย 8 cores (แนะนำ 16 cores หรือมากกว่า)
        \item RAM: อย่างน้อย 32 GB (แนะนำ 64 GB หรือมากกว่า)
        \item Storage: SSD อย่างน้อย 500 GB สำหรับ OS และ Applications
        \item Storage: HDD/SSD เพิ่มเติมสำหรับ VM Storage
        \item Network: การเชื่อมต่อเครือข่ายความเร็วสูงและเสถียร
    \end{mycustomenum2}

    \item \textbf{การติดตั้ง Proxmox VE:}
    \begin{mycustomenum2}
        \item ดาวน์โหลด Proxmox VE ISO จาก \url{https://www.proxmox.com/}
        \item ติดตั้งบน Server Hardware หรือ Virtual Machine
        \item ตั้งค่า Network และ Storage Pool
        \item สร้าง VM Templates สำหรับการใช้งาน
    \end{mycustomenum2}

    \item \textbf{การติดตั้งซอฟต์แวร์พื้นฐาน:}

\begin{figure}[htbp]
\centering
\adjustbox{frame, width=0.9\textwidth}{
    \begin{minipage}{0.85\textwidth}
        \texttt{\footnotesize
        \# Update system\\
        sudo apt update \&\& sudo apt upgrade -y\\
        \\
        \# Install Docker\\
        curl -fsSL https://get.docker.com -o get-docker.sh\\
        sudo sh get-docker.sh\\
        \\
        \# Install Docker Compose\\
        sudo apt install docker-compose-plugin\\
        \\
        \# Install Bun\\
        curl -fsSL https://bun.sh/install | bash
        }
    \end{minipage}
}
\caption{\fontSixTeen{การติดตั้งซอฟต์แวร์พื้นฐานสำหรับ Production}}
\label{figE:SoftwareInstall}
\end{figure}
\end{mycustomenum2}

\section{การตั้งค่า Environment Variables สำหรับ Production}
\hspace*{1.3em}
การตั้งค่า Environment Variables ที่เหมาะสมสำหรับการใช้งานจริง

\vspace{0.5em}
\textbf{ไฟล์ .env สำหรับ Production}
\vspace{0.5em}

\begin{figure}[htbp]
\centering
\adjustbox{frame, width=0.9\textwidth}{
    \begin{minipage}{0.85\textwidth}
        \texttt{\scriptsize
        \# Database Configuration\\
        DB\_HOST=localhost\\
        DB\_PORT=5432\\
        DB\_USER=cloudplatform\\
        DB\_PASSWORD=\{strong\_password\_here\}\\
        DB\_NAME=cloudplatform\_prod\\
        \\
        \# RabbitMQ Configuration\\
        RABBITMQ\_HOST=localhost\\
        RABBITMQ\_PORT=5672\\
        RABBITMQ\_USER=cloudplatform\\
        RABBITMQ\_PASSWORD=\{strong\_password\_here\}\\
        \\
        \# Proxmox Configuration\\
        PROXMOX\_HOST=\{proxmox\_server\_ip\}\\
        PROXMOX\_USERNAME=\{api\_username\}\\
        PROXMOX\_PASSWORD=\{api\_password\}\\
        PROXMOX\_NODE=\{proxmox\_node\_name\}\\
        \\
        \# Application Ports\\
        MIDORI\_PORT=3000\\
        MOMOI\_PORT=3001\\
        \\
        \# Security\\
        JWT\_SECRET=\{very\_strong\_jwt\_secret\}\\
        ENCRYPTION\_KEY=\{32\_char\_encryption\_key\}\\
        \\
        \# Monitoring\\
        JAEGER\_ENDPOINT=http://localhost:14268/api/traces\\
        PROMETHEUS\_PORT=9090\\
        LOKI\_ENDPOINT=http://localhost:3100
        }
    \end{minipage}
}
\caption{\fontSixTeen{ตัวอย่างการตั้งค่า Environment Variables สำหรับ Production}}
\label{figE:EnvProduction}
\end{figure}

\section{การ Deploy ด้วย Docker Compose}
\hspace*{1.3em}
การใช้ Docker Compose สำหรับการ Deploy ระบบแบบ Production

\vspace{0.5em}
\textbf{การสร้าง Production Docker Compose}
\vspace{0.5em}

\begin{mycustomenum2}
    \item สร้างไฟล์ \texttt{docker-compose.prod.yml}

\begin{figure}[htbp]
\centering
\adjustbox{frame, width=0.9\textwidth}{
    \begin{minipage}{0.85\textwidth}
        \texttt{\scriptsize
        version: '3.8'\\
        \\
        services:\\
        \hspace*{1em}midori:\\
        \hspace*{2em}build:\\
        \hspace*{3em}context: .\\
        \hspace*{3em}dockerfile: docker/midori.dockerfile\\
        \hspace*{2em}ports:\\
        \hspace*{3em}- "\$\{MIDORI\_PORT:-3000\}:3000"\\
        \hspace*{2em}environment:\\
        \hspace*{3em}- NODE\_ENV=production\\
        \hspace*{3em}- MOMOI\_API\_URL=http://momoi:3001\\
        \hspace*{2em}depends\_on:\\
        \hspace*{3em}- momoi\\
        \hspace*{2em}restart: unless-stopped\\
        \\
        \hspace*{1em}momoi:\\
        \hspace*{2em}build:\\
        \hspace*{3em}context: .\\
        \hspace*{3em}dockerfile: docker/momoi.dockerfile\\
        \hspace*{2em}ports:\\
        \hspace*{3em}- "\$\{MOMOI\_PORT:-3001\}:3001"\\
        \hspace*{2em}environment:\\
        \hspace*{3em}- NODE\_ENV=production\\
        \hspace*{2em}depends\_on:\\
        \hspace*{3em}- db\\
        \hspace*{3em}- queue\\
        \hspace*{2em}restart: unless-stopped
        }
    \end{minipage}
}
\caption{\fontSixTeen{ตัวอย่าง Docker Compose สำหรับ Production}}
\label{figE:DockerComposeProduction}
\end{figure}

    \item รันคำสั่งสำหรับ Deploy

\begin{figure}[htbp]
\centering
\adjustbox{frame, width=0.9\textwidth}{
    \begin{minipage}{0.85\textwidth}
        \texttt{\footnotesize
        \# Build และ Start Services\\
        docker-compose -f docker-compose.prod.yml up -d --build\\
        \\
        \# ตรวจสอบสถานะ Services\\
        docker-compose -f docker-compose.prod.yml ps\\
        \\
        \# ดู Logs\\
        docker-compose -f docker-compose.prod.yml logs -f
        }
    \end{minipage}
}
\caption{\fontSixTeen{คำสั่งสำหรับ Deploy ด้วย Docker Compose}}
\label{figE:DockerDeployCommands}
\end{figure}
\end{mycustomenum2}

\section{การตั้งค่า Reverse Proxy และ SSL}
\hspace*{1.3em}
การตั้งค่า Nginx เป็น Reverse Proxy และการติดตั้ง SSL Certificate

\vspace{0.5em}
\textbf{การติดตั้งและตั้งค่า Nginx}
\vspace{0.5em}

\begin{mycustomenum2}
    \item ติดตั้ง Nginx

\begin{figure}[htbp]
\centering
\adjustbox{frame, width=0.9\textwidth}{
    \begin{minipage}{0.85\textwidth}
        \texttt{\footnotesize
        sudo apt install nginx\\
        sudo systemctl enable nginx\\
        sudo systemctl start nginx
        }
    \end{minipage}
}
\caption{\fontSixTeen{การติดตั้ง Nginx}}
\label{figE:NginxInstall}
\end{figure}

    \item สร้างไฟล์ Configuration สำหรับ Site

\begin{figure}[htbp]
\centering
\adjustbox{frame, width=0.9\textwidth}{
    \begin{minipage}{0.85\textwidth}
        \texttt{\scriptsize
        server \{\\
        \hspace*{1em}listen 80;\\
        \hspace*{1em}server\_name your-domain.com;\\
        \\
        \hspace*{1em}location / \{\\
        \hspace*{2em}proxy\_pass http://localhost:3000;\\
        \hspace*{2em}proxy\_http\_version 1.1;\\
        \hspace*{2em}proxy\_set\_header Upgrade \$http\_upgrade;\\
        \hspace*{2em}proxy\_set\_header Connection 'upgrade';\\
        \hspace*{2em}proxy\_set\_header Host \$host;\\
        \hspace*{2em}proxy\_set\_header X-Real-IP \$remote\_addr;\\
        \hspace*{2em}proxy\_set\_header X-Forwarded-For \$proxy\_add\_x\_forwarded\_for;\\
        \hspace*{2em}proxy\_set\_header X-Forwarded-Proto \$scheme;\\
        \hspace*{2em}proxy\_cache\_bypass \$http\_upgrade;\\
        \hspace*{1em}\}\\
        \\
        \hspace*{1em}location /api \{\\
        \hspace*{2em}proxy\_pass http://localhost:3001;\\
        \hspace*{2em}proxy\_set\_header Host \$host;\\
        \hspace*{2em}proxy\_set\_header X-Real-IP \$remote\_addr;\\
        \hspace*{2em}proxy\_set\_header X-Forwarded-For \$proxy\_add\_x\_forwarded\_for;\\
        \hspace*{2em}proxy\_set\_header X-Forwarded-Proto \$scheme;\\
        \hspace*{1em}\}\\
        \}
        }
    \end{minipage}
}
\caption{\fontSixTeen{การตั้งค่า Nginx Reverse Proxy}}
\label{figE:NginxConfig}
\end{figure}

    \item ติดตั้ง SSL Certificate ด้วย Let's Encrypt

\begin{figure}[htbp]
\centering
\adjustbox{frame, width=0.9\textwidth}{
    \begin{minipage}{0.85\textwidth}
        \texttt{\footnotesize
        \# ติดตั้ง Certbot\\
        sudo apt install certbot python3-certbot-nginx\\
        \\
        \# ขอ SSL Certificate\\
        sudo certbot --nginx -d your-domain.com\\
        \\
        \# ตั้งค่า Auto-renewal\\
        sudo crontab -e\\
        \# เพิ่มบรรทัด: 0 12 * * * /usr/bin/certbot renew --quiet
        }
    \end{minipage}
}
\caption{\fontSixTeen{การติดตั้ง SSL Certificate}}
\label{figE:SSLInstall}
\end{figure}
\end{mycustomenum2}

\section{การตั้งค่า Database สำหรับ Production}
\hspace*{1.3em}
การเตรียม PostgreSQL Database สำหรับการใช้งานจริง

\vspace{0.5em}
\textbf{การตั้งค่า PostgreSQL}
\vspace{0.5em}

\begin{mycustomenum2}
    \item การตั้งค่า Database User และ Permissions

\begin{figure}[htbp]
\centering
\adjustbox{frame, width=0.9\textwidth}{
    \begin{minipage}{0.85\textwidth}
        \texttt{\footnotesize
        \# เข้าสู่ PostgreSQL\\
        sudo -u postgres psql\\
        \\
        \# สร้าง Database และ User\\
        CREATE DATABASE cloudplatform\_prod;\\
        CREATE USER cloudplatform WITH ENCRYPTED PASSWORD 'strong\_password';\\
        GRANT ALL PRIVILEGES ON DATABASE cloudplatform\_prod TO cloudplatform;\\
        \\
        \# ออกจาก PostgreSQL\\
        \\q
        }
    \end{minipage}
}
\caption{\fontSixTeen{การตั้งค่า PostgreSQL Database}}
\label{figE:PostgreSQLSetup}
\end{figure}

    \item การรัน Database Migration

\begin{figure}[htbp]
\centering
\adjustbox{frame, width=0.9\textwidth}{
    \begin{minipage}{0.85\textwidth}
        \texttt{\footnotesize
        \# เข้าไปใน Momoi directory\\
        cd apps/momoi\\
        \\
        \# รัน Prisma Migration\\
        bunx prisma migrate deploy\\
        \\
        \# Generate Prisma Client\\
        bunx prisma generate
        }
    \end{minipage}
}
\caption{\fontSixTeen{การรัน Database Migration}}
\label{figE:DatabaseMigration}
\end{figure}

    \item การตั้งค่า Database Backup

\begin{figure}[htbp]
\centering
\adjustbox{frame, width=0.9\textwidth}{
    \begin{minipage}{0.85\textwidth}
        \texttt{\footnotesize
        \#!/bin/bash\\
        \# สร้างไฟล์ backup-db.sh\\
        \\
        DB\_NAME="cloudplatform\_prod"\\
        DB\_USER="cloudplatform"\\
        BACKUP\_DIR="/backups/postgresql"\\
        DATE=\$(date +\%Y\%m\%d\_\%H\%M\%S)\\
        \\
        \# สร้าง directory หากไม่มี\\
        mkdir -p \$BACKUP\_DIR\\
        \\
        \# สำรอง Database\\
        pg\_dump -h localhost -U \$DB\_USER -d \$DB\_NAME > \\
        \hspace*{1em}\$BACKUP\_DIR/backup\_\$DATE.sql\\
        \\
        \# ลบไฟล์ backup เก่าที่เก็บเกิน 7 วัน\\
        find \$BACKUP\_DIR -name "backup\_*.sql" -mtime +7 -delete
        }
    \end{minipage}
}
\caption{\fontSixTeen{สคริปต์สำหรับ Database Backup}}
\label{figE:DatabaseBackup}
\end{figure}
\end{mycustomenum2}

\section{การตรวจสอบและบำรุงรักษาระบบ}
\hspace*{1.3em}
การตรวจสอบและบำรุงรักษาระบบเพื่อให้ทำงานได้อย่างเสถียร

\vspace{0.5em}
\textbf{การตรวจสอบสถานะระบบ}
\vspace{0.5em}

\begin{mycustomenum2}
    \item สร้างสคริปต์ตรวจสอบสถานะ Services

\begin{figure}[htbp]
\centering
\adjustbox{frame, width=0.9\textwidth}{
    \begin{minipage}{0.85\textwidth}
        \texttt{\footnotesize
        \#!/bin/bash\\
        \# สร้างไฟล์ health-check.sh\\
        \\
        echo "=== System Health Check ==="\\
        \\
        \# ตรวจสอบ Docker Services\\
        echo "Docker Services Status:"\\
        docker-compose -f docker-compose.prod.yml ps\\
        \\
        \# ตรวจสอบ Nginx\\
        echo "Nginx Status:"\\
        systemctl status nginx --no-pager\\
        \\
        \# ตรวจสอบ Disk Space\\
        echo "Disk Usage:"\\
        df -h\\
        \\
        \# ตรวจสอบ Memory\\
        echo "Memory Usage:"\\
        free -h\\
        \\
        \# ตรวจสอบการเชื่อมต่อ Database\\
        echo "Database Connection:"\\
        pg\_isready -h localhost -p 5432
        }
    \end{minipage}
}
\caption{\fontSixTeen{สคริปต์ตรวจสอบสถานะระบบ}}
\label{figE:HealthCheck}
\end{figure}

    \item ตั้งค่า Cron Job สำหรับการตรวจสอบอัตโนมัติ

\begin{figure}[htbp]
\centering
\adjustbox{frame, width=0.9\textwidth}{
    \begin{minipage}{0.85\textwidth}
        \texttt{\footnotesize
        \# แก้ไข Crontab\\
        crontab -e\\
        \\
        \# เพิ่มงานตรวจสอบทุก 5 นาที\\
        */5 * * * * /path/to/health-check.sh >> /var/log/health-check.log\\
        \\
        \# สำรอง Database ทุกวันเวลา 2:00 น.\\
        0 2 * * * /path/to/backup-db.sh\\
        \\
        \# ทำความสะอาด Docker images เก่าทุกอาทิตย์\\
        0 3 * * 0 docker system prune -f
        }
    \end{minipage}
}
\caption{\fontSixTeen{การตั้งค่า Cron Jobs}}
\label{figE:CronJobs}
\end{figure}
\end{mycustomenum2}

\vspace{0.5em}
\textbf{การอัปเดตระบบ}
\vspace{0.5em}

\begin{mycustomenum2}
    \item การอัปเดตแอปพลิเคชัน

\begin{figure}[htbp]
\centering
\adjustbox{frame, width=0.9\textwidth}{
    \begin{minipage}{0.85\textwidth}
        \texttt{\footnotesize
        \#!/bin/bash\\
        \# สร้างไฟล์ update-app.sh\\
        \\
        echo "Starting application update..."\\
        \\
        \# Pull latest code\\
        git pull origin main\\
        \\
        \# Install dependencies\\
        bun install\\
        \\
        \# Build และ restart services\\
        docker-compose -f docker-compose.prod.yml down\\
        docker-compose -f docker-compose.prod.yml up -d --build\\
        \\
        \# ตรวจสอบสถานะหลัง update\\
        sleep 30\\
        docker-compose -f docker-compose.prod.yml ps\\
        \\
        echo "Update completed!"
        }
    \end{minipage}
}
\caption{\fontSixTeen{สคริปต์สำหรับอัปเดตแอปพลิเคชัน}}
\label{figE:UpdateApp}
\end{figure}
\end{mycustomenum2}

\section{การแก้ไขปัญหาใน Production}
\hspace*{1.3em}
วิธีการแก้ไขปัญหาที่อาจเกิดขึ้นในการใช้งานจริง

\vspace{0.5em}
\textbf{ปัญหาที่พบบ่อยและวิธีแก้ไข}
\vspace{0.5em}

\begin{mycustomenum2}
    \item \textbf{Service ไม่ตอบสนอง}
    \begin{mycustomenum2}
        \item ตรวจสอบ Logs: \texttt{docker-compose logs [service\_name]}
        \item รีสตาร์ท Service: \texttt{docker-compose restart [service\_name]}
        \item ตรวจสอบ Resource Usage: \texttt{docker stats}
    \end{mycustomenum2}

    \item \textbf{Database Connection Error}
    \begin{mycustomenum2}
        \item ตรวจสอบ PostgreSQL Service: \texttt{systemctl status postgresql}
        \item ตรวจสอบ Connection Pool: ดู Logs ของ Momoi Service
        \item รีสตาร์ท Database: \texttt{sudo systemctl restart postgresql}
    \end{mycustomenum2}

    \item \textbf{Out of Disk Space}
    \begin{mycustomenum2}
        \item ทำความสะอาด Docker: \texttt{docker system prune -a}
        \item ลบ Log files เก่า: \texttt{find /var/log -name "*.log" -mtime +30 -delete}
        \item ตรวจสอบ VM Storage ใน Proxmox
    \end{mycustomenum2}

    \item \textbf{SSL Certificate หมดอายุ}
    \begin{mycustomenum2}
        \item ต่ออายุ Certificate: \texttt{sudo certbot renew}
        \item รีโหลด Nginx: \texttt{sudo nginx -s reload}
        \item ตรวจสอบวันหมดอายุ: \texttt{sudo certbot certificates}
    \end{mycustomenum2}
\end{mycustomenum2}